\documentclass[aspectratio=169,11pt]{beamer}
% Shared CMPSC 480 Beamer preamble.
\usepackage[T1]{fontenc}
\usepackage[utf8]{inputenc}
\usepackage{lmodern}
\usepackage{amsmath,amssymb,mathtools}
\usepackage{booktabs,array,tabularx}
\usepackage{xcolor}
\usepackage{listings}
\usepackage{tikz}
\usetikzlibrary{arrows.meta,positioning,shapes.geometric}

% Ignore only negligible TeX box diagnostics that are below visual tolerance.
% Larger layout defects still appear in the build log and in rendered QA.
\vfuzz=3pt
\hbadness=2000

\definecolor{courseink}{HTML}{111111}
\definecolor{coursegray}{HTML}{EDEDED}
\definecolor{courserule}{HTML}{B8BCC4}
\definecolor{courseblue}{HTML}{3D8DFF}
\definecolor{courselightblue}{HTML}{D0EDFA}
\definecolor{coursemuted}{HTML}{5C6470}
\definecolor{coursegreen}{HTML}{0B7A53}
\definecolor{coursered}{HTML}{B3261E}

\usetheme{default}
\usefonttheme{professionalfonts}
\setbeamertemplate{navigation symbols}{}
\setbeamersize{text margin left=0.65cm,text margin right=0.65cm}

\setbeamercolor{normal text}{fg=courseink,bg=white}
\setbeamercolor{frametitle}{fg=courseink,bg=white}
\setbeamercolor{title}{fg=courseink,bg=white}
\setbeamercolor{subtitle}{fg=coursemuted,bg=white}
\setbeamercolor{structure}{fg=courseblue}
\setbeamercolor{alerted text}{fg=coursered}
\setbeamercolor{block title}{fg=courseink,bg=courselightblue}
\setbeamercolor{block body}{fg=courseink,bg=coursegray}
\setbeamercolor{example text}{fg=coursegreen}

\setbeamerfont{title}{size=\fontsize{25}{29}\selectfont,series=\bfseries}
\setbeamerfont{subtitle}{size=\fontsize{13}{16}\selectfont}
\setbeamerfont{frametitle}{size=\fontsize{19}{22}\selectfont,series=\bfseries}
\setbeamerfont{normal text}{size=\fontsize{11}{14}\selectfont}
\setbeamerfont{footline}{size=\fontsize{7.5}{9}\selectfont}

\setbeamertemplate{frametitle}{%
  \nointerlineskip
  % Use content-driven height so a long assessment title can wrap without
  % being clipped by a fixed-height title bar.
  \begin{beamercolorbox}[wd=\paperwidth,sep=0.17cm,leftskip=0.48cm,rightskip=0.48cm]{frametitle}
    \insertframetitle
  \end{beamercolorbox}
  \vspace{-0.08cm}
  {\color{courserule}\rule{\textwidth}{0.35pt}}
}

\setbeamertemplate{footline}{%
  \leavevmode
  \begin{beamercolorbox}[wd=\paperwidth,ht=2.6ex,dp=1.1ex,leftskip=.65cm,rightskip=.65cm]{author in head/foot}
    \color{coursemuted}\insertshorttitle\hfill\insertframenumber/\inserttotalframenumber
  \end{beamercolorbox}
}

\setbeamertemplate{itemize item}{\color{courseblue}\small$\blacktriangleright$}
\setbeamertemplate{itemize subitem}{\color{courseblue}\scriptsize$\bullet$}
\setbeamertemplate{enumerate item}{\color{courseblue}\insertenumlabel.}

\lstdefinestyle{coursepython}{
  language=Python,
  basicstyle=\ttfamily\fontsize{8.5}{10}\selectfont,
  keywordstyle=\color{courseblue}\bfseries,
  commentstyle=\color{coursemuted}\itshape,
  stringstyle=\color{coursegreen},
  showstringspaces=false,
  columns=fullflexible,
  keepspaces=true,
  frame=single,
  rulecolor=\color{courserule},
  backgroundcolor=\color{coursegray},
  breaklines=true,
  tabsize=4,
  xleftmargin=0.15cm,
  xrightmargin=0.15cm,
  framexleftmargin=0.15cm,
  framexrightmargin=0.15cm
}
\lstset{style=coursepython}

\newcommand{\points}[1]{\hfill{\color{courseblue}\textbf{[#1 points]}}}
\newcommand{\deliverable}[1]{\textbf{\color{courseblue}#1}}
\newcommand{\code}[1]{\texttt{#1}}
\newcommand{\keyidea}[1]{%
  \begin{beamercolorbox}[rounded=false,sep=0.55em,wd=\linewidth]{block body}
    \textbf{Key idea.} #1
  \end{beamercolorbox}
}
\newcommand{\answerlabel}{\textbf{\color{coursegreen}Answer.}\ }
\newcommand{\gradinglabel}{\textbf{\color{courseblue}Grading guidance.}\ }

\newenvironment{questionblock}[1]
  {\begin{block}{#1}}
  {\end{block}}

\newenvironment{solutionblock}[1]
  {\begin{block}{#1}}
  {\end{block}}

\AtBeginSection[]{
  \begin{frame}
    \vfill
    {\usebeamerfont{title}\color{courseink}\insertsectionhead\par}
    \vspace{0.4cm}
    {\color{courseblue}\rule{0.32\paperwidth}{3pt}}
    \vfill
  \end{frame}
}


% CMPSC480_TOTAL_POINTS: 10
% CMPSC480_ITEM_IDS: 1,2,3,4
% CMPSC480_SUBITEM_POINTS: 1=2,2=5,3=1,4=2
\title[CMPSC 480 - Week 1 Lab]{Week 1 Code Lab}
\subtitle{Environment Setup and Your First Agent/Problem Skeleton}
\author{CMPSC 480 - Student Question Deck}
\date{}

\begin{document}

\begin{frame}
  \titlepage
\end{frame}

\begin{frame}{Assignment overview}
  \begin{columns}[T,onlytextwidth]
    \begin{column}{0.58\textwidth}
      \Large
      Build a verified Python development environment and a reusable
      agent/problem abstraction for a small gridworld.

      \vspace{0.45cm}
      \normalsize
      This is a tooling and representation lab. You will not implement a
      search algorithm yet.
    \end{column}
    \begin{column}{0.37\textwidth}
      \begin{block}{Expected effort}
        2-3 hours
      \end{block}
      \begin{block}{Points}
        10 total
      \end{block}
      \begin{block}{Submission}
        Git repository archive or repository link
      \end{block}
    \end{column}
  \end{columns}
\end{frame}

\begin{frame}{Learning goals}
  \begin{itemize}
    \item Verify a reproducible Python, NumPy, Jupyter, and Git environment.
    \item Map a formal search problem onto a concrete Python interface.
    \item Represent states as immutable, hashable values.
    \item Define legal actions, state transitions, a goal test, and step costs.
    \item Use validation, type hints, docstrings, and unit tests from the start.
    \item Create abstractions that can be extended in Week 2 without redesign.
  \end{itemize}
\end{frame}

\begin{frame}{Required repository structure}
  \begin{block}{Submit this structure}
    \ttfamily
    week01-agent-skeleton/\par
    \hspace{0.45cm}README.md\par
    \hspace{0.45cm}verify\_environment.py\par
    \hspace{0.45cm}src/\par
    \hspace{0.9cm}\_\_init\_\_.py\par
    \hspace{0.9cm}agent.py\par
    \hspace{0.9cm}problem.py\par
    \hspace{0.45cm}tests/\par
    \hspace{0.9cm}test\_problem.py
  \end{block}
  \vspace{0.25cm}
  \begin{itemize}
    \item The repository must contain at least one meaningful Git commit.
    \item The README must state the Python version and exact test command.
    \item Do not commit virtual environments, cache folders, or credentials.
  \end{itemize}
\end{frame}

\begin{frame}{Checklist - Environment and structure}
  \begin{enumerate}
    \item \textbf{Environment report:} \texttt{verify\_environment.py} reports
      Python, NumPy, Jupyter, and Git status.
    \item \textbf{Version enforcement:} the script exits unsuccessfully when
      Python is older than 3.10 or a required package is unavailable.
    \item \textbf{Git evidence:} the script confirms that it runs inside a Git
      repository with at least one commit.
    \item \textbf{State contract:} every grid state is immutable, hashable, and
      comparable by value.
  \end{enumerate}
\end{frame}

\begin{frame}{Checklist - Behavior and quality}
  \begin{enumerate}
    \setcounter{enumi}{4}
    \item \textbf{Problem interface:} implement
      \texttt{initial\_state()}, \texttt{actions(state)},
      \texttt{result(state, action)}, \texttt{goal\_test(state)}, and
      \texttt{step\_cost(state, action, next\_state)}.
    \item \textbf{Defensive behavior:} reject malformed states and illegal
      actions with precise exceptions; never mutate a caller-provided state.
    \item \textbf{Agent interface:} define an extensible \texttt{Agent} base
      class with an \texttt{act(percept)} contract.
    \item \textbf{Verification and documentation:} provide type hints,
      public-method docstrings, and at least five focused unit tests.
  \end{enumerate}
\end{frame}

\begin{frame}{The toy domain}
  \begin{columns}[T,onlytextwidth]
    \begin{column}{0.46\textwidth}
      \begin{block}{Gridworld}
        Use a rectangular grid with:
        \begin{itemize}
          \item integer coordinates \((row, column)\),
          \item four actions: up, down, left, right,
          \item optional blocked cells,
          \item one initial cell and one goal cell,
          \item unit cost for every legal move.
        \end{itemize}
      \end{block}
    \end{column}
    \begin{column}{0.49\textwidth}
      \begin{block}{Required conventions}
        \begin{itemize}
          \item The upper-left cell is \((0,0)\).
          \item Rows increase downward.
          \item Columns increase to the right.
          \item \texttt{actions(state)} returns an empty list when no move is legal.
          \item The Week 1 code contains no search loop.
        \end{itemize}
      \end{block}
    \end{column}
  \end{columns}
\end{frame}

\begin{frame}{Task 1 - Environment verification (2 points)}
  Create \texttt{verify\_environment.py}. A successful run must:
  \begin{enumerate}
    \item print the complete Python version;
    \item reject Python versions earlier than 3.10;
    \item print installed NumPy and Jupyter versions;
    \item confirm that the current directory is inside a Git work tree;
    \item report a positive commit count; and
    \item exit with status code zero only when every check passes.
  \end{enumerate}
  \vspace{0.25cm}
  \keyidea{The output must be machine-checkable: use stable labels and return
  a nonzero process status for a failed requirement.}
\end{frame}

\begin{frame}{Task 2 - Problem abstraction (5 points)}
  Implement a concrete \texttt{GridProblem} that owns the domain configuration
  and exposes the required search-problem operations.

  \vspace{0.3cm}
  \begin{itemize}
    \item Validate dimensions, initial state, goal state, and blocked cells.
    \item Return legal actions in a documented deterministic order.
    \item Return a new immutable state from \texttt{result}.
    \item Reject an action that is not legal in the supplied state.
    \item Make the goal test return a real Boolean value.
    \item Return a positive numeric step cost only for a legal transition.
  \end{itemize}
\end{frame}

\begin{frame}{Task 3 - Agent abstraction (1 point)}
  Create an \texttt{Agent} base class with this public contract:

  \vspace{0.35cm}
  \begin{block}{act(percept)}
    Accept the agent's current percept and return the next action selected by a
    concrete subclass.
  \end{block}

  \begin{itemize}
    \item The base implementation must not pretend to choose an action.
    \item Document the expected percept and action types.
    \item Design the interface so Week 2 search agents can subclass it.
    \item Do not implement BFS, DFS, UCS, Greedy Best-First Search, or A-star.
  \end{itemize}
\end{frame}

\begin{frame}{Task 4 - Contract tests and repository quality (2 points)}
  Provide at least five focused tests. Your suite must include:
  \begin{itemize}
    \item a known legal transition;
    \item a corner or blocked-cell action set;
    \item an illegal-action exception;
    \item a goal-test success and failure;
    \item evidence that a state can be used as a set member or dictionary key;
    \item a domain with no legal actions;
    \item at least one malformed-input rejection; and
    \item a reproducible README with setup and test commands, plus a clean,
      meaningful Git commit.
  \end{itemize}

  \vspace{0.25cm}
  You may combine closely related assertions in one test, but each failure must
  remain easy to diagnose.
\end{frame}

\begin{frame}[fragile]{Starter code - state and action types}
\begin{lstlisting}
from __future__ import annotations

from enum import Enum
from typing import TypeAlias

State: TypeAlias = tuple[int, int]


class Action(Enum):
    UP = (-1, 0)
    DOWN = (1, 0)
    LEFT = (0, -1)
    RIGHT = (0, 1)
\end{lstlisting}
\end{frame}

\begin{frame}[fragile]{Starter code - problem interface}
\begin{lstlisting}
class GridProblem:
    def initial_state(self) -> State:
        raise NotImplementedError

    def actions(self, state: State) -> list[Action]:
        raise NotImplementedError

    def result(self, state: State, action: Action) -> State:
        raise NotImplementedError

    def goal_test(self, state: State) -> bool:
        raise NotImplementedError

    def step_cost(
        self, state: State, action: Action, next_state: State
    ) -> float:
        raise NotImplementedError
\end{lstlisting}
\end{frame}

\begin{frame}[fragile]{Starter code - agent interface}
\begin{lstlisting}
from abc import ABC, abstractmethod
from typing import Generic, TypeVar

PerceptT = TypeVar("PerceptT")
ActionT = TypeVar("ActionT")


class Agent(ABC, Generic[PerceptT, ActionT]):
    @abstractmethod
    def act(self, percept: PerceptT) -> ActionT:
        """Return the next action for the supplied percept."""
        raise NotImplementedError
\end{lstlisting}
\end{frame}

\begin{frame}{Edge cases and defensive-programming rules}
  \begin{itemize}
    \item A state outside the grid is malformed even if its tuple shape is valid.
    \item A blocked cell is not a valid current, initial, goal, or next state.
    \item \texttt{result} must reject an action not returned by
      \texttt{actions(state)}.
    \item \texttt{goal\_test} must not silently accept malformed input.
    \item A one-cell grid may have an empty action list and may begin at the goal.
    \item State values must remain safe as dictionary keys after every operation.
  \end{itemize}
\end{frame}

\begin{frame}{Submission instructions}
  \begin{enumerate}
    \item Run \texttt{python verify\_environment.py} from the repository root.
    \item Run the complete test suite using the command documented in README.
    \item Confirm that generated caches and virtual environments are ignored.
    \item Commit all required source and test files with a meaningful message.
    \item Submit the repository link or archive through the course system.
  \end{enumerate}

  \vspace{0.3cm}
  \begin{alertblock}{Do not submit}
    API keys, access tokens, editor settings, virtual environments, or code for
    any search algorithm.
  \end{alertblock}
\end{frame}

\begin{frame}{Grading rubric - Correctness (8 points)}
  \scriptsize
  \begin{tabularx}{\textwidth}{>{\bfseries}p{0.27\textwidth} c X}
    \toprule
    Requirement & Points & Evidence \\
    \midrule
    Environment verification & 2 & Versions, Git status, commit count, and
      meaningful failure status \\
    Immutable state model & 1 & Hashable value semantics with no caller-visible mutation \\
    Legal actions and transitions & 2 & Correct boundaries, blocked cells, and
      deterministic actions \\
    Goal test and step cost & 1 & Correct Boolean result and legal unit-cost transition \\
    Defensive validation & 1 & Precise rejection of malformed states and illegal actions \\
    Agent abstraction & 1 & Extensible, documented \texttt{act} contract \\
    \bottomrule
  \end{tabularx}
\end{frame}

\begin{frame}{Rubric - Verification and quality (2 points)}
  \scriptsize
  \begin{tabularx}{\textwidth}{>{\bfseries}p{0.27\textwidth} c X}
    \toprule
    Requirement & Points & Evidence \\
    \midrule
    Focused unit tests & 1 & Required normal, boundary, illegal, no-action,
      hashability, and malformed cases \\
    Documentation and repository quality & 1 & Type hints, docstrings, README,
      clean structure, and meaningful commit \\
    \midrule
    \textbf{Total} & \textbf{10} & \\
    \bottomrule
  \end{tabularx}
\end{frame}

\begin{frame}{Reflection prompt}
  In 150-250 words:

  \vspace{0.35cm}
  \begin{enumerate}
    \item Map the formal search-problem components - initial state, actions,
      transition model, goal test, and path cost - onto your class design.
    \item Identify one ambiguity you resolved while defining the gridworld contract.
    \item Explain why your state representation will work with an explored set
      in next week's graph-search implementation.
  \end{enumerate}

  \vspace{0.3cm}
  Submit the response in \texttt{README.md}.
\end{frame}

\end{document}
